\documentclass[UTF8]{ctexart}
\usepackage{amsmath,amssymb}
\usepackage[a4paper,margin=2.5cm]{geometry}
\title{空间直角坐标，平面和直线}
\date{}

\begin{document}
\maketitle

\subsection*{§1\quad 空间直角坐标系}

\begin{itemize}
\item 在空间中，固定点 $O$，过 $O$ 点作三条两两垂直的数轴
$OX,OY,OZ$，此时有空间直角坐标系，记作 $[O;X,Y,Z]$；
\item 称 $O$ 是坐标原点，$OX,OY,OZ$ 为坐标轴；
\item $OX$ 与 $OY$ 形成的平面为 $XOY$ 面，以此类推，
它们是“坐标面”；
\item $XOY$ 面将空间分成两部分，$Z$ 轴所指的部分称为正面，
另外一面称为反面；
\item 任给空间内一点 $P$，它到 $YOZ$ 面、$XOZ$ 面及 $XOY$ 面的
“距离”（点 $P$ 在正面为正，在反面为负）分别为 $x,y,z$，
称三元实数对 $(x,y,z)$ 是点 $P$ 在 $[O;X,Y,Z]$ 下的坐标；
\item 显然 $|OP|=\sqrt{x^2+y^2+z^2}$；
\item 右手系与左手系。
\end{itemize}

\textbf{eg.}\quad 对于空间内的点 $P_1(x_1,y_1,z_1)$、
$P_2(x_2,y_2,z_2)$，如何求 $|P_1P_2|$？

过点 $P_1$ 作三条分别平行于 $x$ 轴、$y$ 轴、$z$ 轴的数轴，
记为 $x'$ 轴、$y'$ 轴、$z'$ 轴，此时有新的坐标系
$[P_1;X',Y',Z']$，$P_1$ 在 $[P_1;X',Y',Z']$ 中的坐标为
$(0,0,0)$，$P_2$ 在 $[P_1;X',Y',Z']$ 中的坐标为
$(x_2',y_2',z_2')$，于是有
\[
\begin{cases}
x_2'=x_2-x_1,\\
y_2'=y_2-y_1,\\
z_2'=z_2-z_1,
\end{cases}
\]
这是\textbf{坐标平移公式}。得到
\[
|P_1P_2|=\sqrt{(x_2-x_1)^2+(y_2-y_1)^2+(z_2-z_1)^2},
\]
即两点距离公式。

\subsection*{§2\quad 方向的表示}

\subsubsection*{1.\ 方向}

由起点和终点确定的射线所表示。

\begin{itemize}
\item 如果起点是原点，终点是 $P(x,y,z)$，用三元实数对
$(x,y,z)$ 表示该方向：当 $\lambda>0$ 时，
$(\lambda x,\lambda y,\lambda z)$ 也表示该方向；当
$\lambda<0$ 时，$(\lambda x,\lambda y,\lambda z)$ 表示
$OP$ 方向的反方向；
\item 如果起点是 $P_1(x_1,y_1,z_1)$，终点
$P_2(x_2,y_2,z_2)$，用
$(x_2-x_1,\ y_2-y_1,\ z_2-z_1)$ 表示方向 $P_1P_2$。
\end{itemize}

\subsubsection*{2.\ 方向角}

方向 $OP$ 分别与 $X$ 轴、$Y$ 轴、$Z$ 轴的夹角，记为
$\alpha,\beta,\gamma$。

\begin{itemize}
\item 若 $P(x,y,z)$，则 $\cos\alpha=\dfrac{x}{|OP|}$，
$\beta,\gamma$ 也同理。显然
$\cos^2\alpha+\cos^2\beta+\cos^2\gamma=1$；
\item 称 $\cos\alpha,\cos\beta,\cos\gamma$ 是方向 $OP$
的\textbf{方向余弦}；
\item 方向余弦是方向的单位化，它在单位球面上。
\end{itemize}

\subsubsection*{3.\ 方向的夹角}

若两个方向的方向余弦分别为 $(l_1,m_1,n_1)$ 及
$(l_2,m_2,n_2)$，它们的夹角记为 $\theta$，则一方面
\[
|P_1P_2|^2=(l_2-l_1)^2+(m_2-m_1)^2+(n_2-n_1)^2 ;
\]
另一方面，我们可以用余弦定理
\[
|P_1P_2|^2=|OP_1|^2+|OP_2|^2-2|OP_1||OP_2|\cos\theta .
\]
于是
\[
1+1-2l_1l_2-2m_1m_2-2n_1n_2=1+1-2\cos\theta
\ \Longrightarrow\
\cos\theta=l_1l_2+m_1m_2+n_1n_2 .
\]
当 $\theta=\dfrac\pi2$ 时，称上述两方向垂直；若 $\theta=0$，
则称为平行。

\subsection*{§3\quad 空间中的平面}

\begin{itemize}
\item \textbf{平面}\quad 到定点的方向与给定方向垂直的点的集合；
\item 给定的方向称为平面的\textbf{法向}。不难看出，
一个平面有两个法向，方向相反。
\end{itemize}

已知点 $O$ 在平面 $\pi$ 的垂足为 $N$，且 $ON$ 的方向余弦为
$(l,m,n)$，$|ON|=p$。任取点 $P(x,y,z)\in\pi$，
则方向 $NP$ 与方向 $ON$ 垂直，于是
\[
(x-pl)l+(y-pm)m+(z-pn)n=0
\ \Longleftrightarrow\
lx+my+nz-p=0 .
\]

\textbf{注：}

\begin{enumerate}
\item[①] 若 $p=0$（即 $O\in\pi$），则法式方程不唯一，
它差一个符号；
\item[②] 若 $p>0$，则法式方程唯一，其几何意义是：
与半径为 $p$ 的球面相切于方向是 $(l,m,n)$ 的点的平面；
\item[③] 法式方程的特点：\\
1. 一次项系数平方和为 $1$；\quad 2. 常数项 $\leq0$。
\end{enumerate}

\subsubsection*{特殊的平面}

\begin{enumerate}
\item $Ax+By+D=0$ 表示与 $z$ 轴平行的平面。
\item 若
\[
\begin{cases}
\pi_1:A_1x+B_1y+C_1z+D_1=0,\\
\pi_2:A_2x+B_2y+C_2z+D_2=0,
\end{cases}
\]
则有
\[
\pi_1\parallel\pi_2
\Longleftrightarrow
\frac{A_1}{A_2}=\frac{B_1}{B_2}=\frac{C_1}{C_2},
\qquad
\pi_1\perp\pi_2
\Longleftrightarrow
A_1A_2+B_1B_2+C_1C_2=0 .
\]
\item 如果
\[
\begin{cases}
\pi_1:A_1x+B_1y+C_1z+D_1=0,\\
\pi_2:A_2x+B_2y+C_2z+D_2=0,
\end{cases}
\]
则 $\pi_1$ 到 $\pi_2$ 的距离是
\[
\frac{|D_1-D_2|}{\sqrt{A^2+B^2+C^2}} .
\]
\item 如果
\[
\begin{cases}
\pi_1:A_1x+B_1y+C_1z+D_1=0,\\
\pi_2:A_1x+B_1y+D_1=0,
\end{cases}
\qquad\text{且}\ C\neq0,
\]
那么 $\pi_2$ 是过 $\pi_1$ 与 $XOY$ 面的交线且与 $Z$ 轴平行的平面。
\end{enumerate}

法式方程是一个特殊的三元一次方程，它表示一个平面。
那么，一个一般的三元一次方程
\[
Ax+By+Cz+D=0
\]
表示什么图形？
\[
\frac{\pm1}{\sqrt{A^2+B^2+C^2}}(Ax+By+Cz+D)=0
\]
是一个法式方程，它表示一个平面。称其为平面的\textbf{一般方程}。

\subsection*{§4\quad 空间中的直线}

\begin{itemize}
\item \textbf{直线}\quad 到定点的方向与给定方向平行的点的轨迹。
\end{itemize}

记直线 $L$ 过定点 $N(x_0,y_0,z_0)$，直线 $L$ 与方向
$(l,m,n)$ 平行，取点 $P(x,y,z)\in L$，则方向 $(l,m,n)$
与方向 $NP$ 平行，于是
\[
\exists t\in\mathbb R,\ \text{s.t.}\
\begin{cases}
x-x_0=tl,\\
y-y_0=tm,\\
z-z_0=tn,
\end{cases}
\qquad\text{其中}\ t\in\mathbb R\ \text{是参数} .
\]
称上述是直线 $L$ 的\textbf{参数方程}。消去参数 $t$，得
\[
\frac{x-x_0}{l}=\frac{y-y_0}{m}=\frac{z-z_0}{n},
\]
若分母为零，则分子也为零。称其为直线 $L$ 的\textbf{点向方程}。

可以把点向方程视作两个方程组成的三元一次方程组。

\textbf{Q：}\quad 一般的，由两个方程组成的三元一次方程组
\[
\begin{cases}
A_1x+B_1y+C_1z+D_1=0,\\
A_2x+B_2y+C_2z+D_2=0,
\end{cases}
\]
表示什么图形？从几何上看，它表示两个平面的交线（如果有的话）；
从代数上看，它可以化成参数方程／点向方程，它是一条直线。
上式称为直线的\textbf{一般方程}。

注意到，若 $L$ 与 $XOY$ 面不平行
$\Longleftrightarrow A_1B_2\neq A_2B_1$，以后可以用外积证明。

\subsubsection*{直线的投影}

有直线 $L$
\[
\begin{cases}
A_1x+B_1y+C_1z+D_1=0,\\
A_2x+B_2y+C_2z+D_2=0,
\end{cases}
\]
上两式联立，消去 $z$，得到方程
\[
A_3x+B_3y+D_3=0,
\]
这是过直线 $L$ 且与 $z$ 轴平行的平面。则方程
\[
\begin{cases}
A_3x+B_3y+D_3=0,\\
z=0,
\end{cases}
\]
的解即为直线在 $XOY$ 平面上的投影。

事实上，就是把点向方程擦掉含 $z$ 的那一部分。

\end{document}
